\documentclass[11pt]{article}
\usepackage[T1]{fontenc}
\usepackage[utf8]{inputenc}
\usepackage{lmodern}
\usepackage[margin=1in]{geometry}
\usepackage{amsmath,amssymb,amsthm,booktabs,array}
\usepackage{microtype}
\usepackage[numbers]{natbib}
\usepackage[hidelinks]{hyperref}
\usepackage{url}
\usepackage{fancyvrb}
\hypersetup{pdftitle={Scalar Hall Exclusion in Dimensions 11 through 15: A Partial Lean Formalization},pdfsubject={A technical note on a verified scalar theorem},pdfauthor={Ryosuke Shiroshita}}
\newtheorem{theorem}{Theorem}
\newtheorem{lemma}[theorem]{Lemma}
\newtheorem{corollary}[theorem]{Corollary}
\newcommand{\Q}{\mathbb{Q}}
\newcommand{\R}{\mathbb{R}}
\newcommand{\N}{\mathbb{N}}
\newcommand{\code}[1]{\texttt{\small #1}}
\setlength{\emergencystretch}{1.5em}
\title{Scalar Hall Exclusion in Dimensions 11 through 15:\\A Partial Lean Formalization}
\author{Ryosuke Shiroshita\\[2pt]
\small\href{mailto:Ryosuke.Shiroshita@exactory.ai}{\texttt{Ryosuke.Shiroshita@exactory.ai}}}
\date{}

\begin{document}
\maketitle
\begin{abstract}
We formalize a finite scalar layer of an existing computer-assisted argument
for Dittert's conjecture. For dimensions 11 through 15, the artifact checks
direct rational formulas for every ordered proper cut and one normalized
marginal cut, then derives a strict polynomial inequality for every real
parameter $u\leq 1$. Its corollary excludes simultaneous permanent and
deficit lower bounds even under a relaxed scalar budget. Lean verifies
340 normalized cases and proves that they cover the entire specified
nonfull scalar domain. Independent execution and statement review confirm
that the positivity, exclusion, and finite-certificate declarations depend
only on \code{propext}, \code{Classical.choice}, and \code{Quot.sound}.
A separate appendix verifies two literal near-zero polynomial
certificates in dimensions 8 and 9 on specified closed intervals;
it does not establish full scalar coverage in either dimension.
The contribution is an attributed partial formalization and reproducible
verification artifact. The matrix reduction, its analytic and structural
inputs, and the full Dittert theorem are outside the formalization.
\end{abstract}

\section{Source argument and contribution}

For a nonnegative real $n\times n$ matrix $A$ of total entry sum $n$, let
$r_i$ and $c_j$ be its row and column sums. Dittert's functional is
\[
 \Phi(A)=\prod_{i=1}^{n}r_i+\prod_{j=1}^{n}c_j-\operatorname{per}(A).
\]
The conjecture asserts $\Phi(A)\leq 2-n!/n^n$, with equality precisely at
the matrix whose entries are $1/n$. Cheon and Wanless establish structural
restrictions on potential maximizers, including exclusions associated
with partial decomposability and rectangular zero sets
\cite{cheon2012some}. Pang's scaling argument provides related
large-dimension context \cite{pang2026proof}.

This note concerns a scalar calculation in Nascimento's public working
proof, at a fixed source revision.\footnote{\raggedright Pedro Paulo Marques do
Nascimento, \href{https://github.com/pedromnasc/dittert-conjecture-proof/tree/894066bbaa715138c98bf3cb7c6fdb4f39a37701}{\texttt{pedromnasc/dittert-conjecture-proof}},
\par\noindent Source commit: \nolinkurl{894066bbaa715138c98bf3cb7c6fdb4f39a37701}.}
That repository already presents an assembled working-proof chain and
states the review boundaries of its components. Its Hall argument
associates a possible boundary maximizer with cut sizes and a scaling
parameter, bounds its permanent from below, and bounds its deficit from
below. A strict scalar inequality makes those bounds incompatible with
the available budget. The finite arithmetic and this final scalar
implication are the objects formalized here.

The source mathematics is not a new discovery of this artifact.
In particular, the note does not claim a new mathematical theorem or
the first formalization of Dittert's conjecture. Li, Xiong, and Yang
already report a dimension-four proof based on an exact sum-of-squares
certificate and a Lean formalization \cite{li2026proof}. Our scope is a
different, explicitly bounded scalar layer of the public Hall argument.
Neither the name of a scalar variable nor a successful arithmetic
check establishes its relationship to an arbitrary matrix.

The formalization uses Lean~4 \cite{moura2021lean} and Mathlib
\cite{community2020lean}. The implementation combines direct
dimension-dependent formulas, exhaustive coverage, rational-to-real
transport, a completed-square argument, Bernoulli comparison, and the
relaxed-budget corollary. The accompanying evidence records
both the exact theorem and what remains outside it.

\section{Exact scalar statement}

All factorials and powers below have natural indices. The constants are
first defined in $\Q$ and then mapped into $\R$ by the canonical embedding.
For the positive indices used below, define
\begin{equation}\label{eq:gamma}
 \gamma_k=\frac{k!}{k^k}.
\end{equation}
Fix $n,a,b\in\N$ satisfying
\begin{equation}\label{eq:domain}
 11\leq n\leq 15,\qquad 0<a+b<n.
\end{equation}
Put $p=n-a-b$. The scalar floor and curvature are
\begin{equation}\label{eq:floor}
 M_{n,a,b}=
 \begin{cases}
 \displaystyle\frac{\gamma_{n-a}\gamma_{n-b}}{\gamma_{n-a-b}},
       & a>0,\ b>0,\\[6pt]
 \displaystyle(n-2)!\left(\frac{n-2}{(n-1)^2}\right)^{n-2},
       & a=0\ \text{or}\ b=0,
 \end{cases}
\end{equation}
and
\begin{equation}\label{eq:curvature}
 C_{n,a,b}=
 \begin{cases}
 \displaystyle\frac{7n(n-a-b)^2}
 {16\{a(n-a)+b(n-b)\}},& a>0,\ b>0,\\[6pt]
 \displaystyle\frac{7n}{16(n-1)}
       \bigl(1+(1-\gamma_n)^2\bigr),& a=0\ \text{or}\ b=0.
 \end{cases}
\end{equation}
The domain ensures that $p$, $n-a$, and $n-b$ are positive. Every
denominator displayed in the proper branch is therefore positive;
the marginal branch has $n-1>0$. Natural subtraction in the Lean
definitions agrees with ordinary subtraction at all these indices.
Although the definitions are total Lean functions, no assertion here
uses out-of-domain values to avoid a denominator obligation.

\begin{theorem}[Scalar positivity]\label{thm:positive}
For every $n,a,b\in\N$ satisfying~\eqref{eq:domain} and every $u\in\R$
with $u\leq 1$,
\begin{equation}\label{eq:positive}
 C_{n,a,b}u^2+M_{n,a,b}(1-u)^n-\gamma_n>0.
\end{equation}
\end{theorem}

\begin{corollary}[Relaxed-budget exclusion]\label{cor:exclusion}
Under the same natural-size restrictions, there are no real $u,P,\delta$
with $u\leq 1$ satisfying all three inequalities
\begin{equation}\label{eq:budget}
 P\geq M_{n,a,b}(1-u)^n,\qquad
 \delta\geq C_{n,a,b}u^2,\qquad P+\delta\leq\gamma_n.
\end{equation}
\end{corollary}

\begin{proof}
The first two inequalities imply
$P+\delta-\gamma_n\geq M_{n,a,b}(1-u)^n+C_{n,a,b}u^2-\gamma_n>0$
by Theorem~\ref{thm:positive}, contradicting the last inequality.
\end{proof}

There is no lower bound on $u$ in either statement. The polynomial
result includes both endpoints of $[0,1]$ and every negative $u$.
This larger scalar domain does not assert that the matrix construction
has a corresponding interpretation there. Likewise, $P$ and $\delta$
are arbitrary real variables in Corollary~\ref{cor:exclusion}; no
definition identifying them with matrix quantities occurs in the theorem.

\section{Finite criterion and real-variable proof}

The finite calculation proves a strict discriminant condition. Writing
$M=M_{n,a,b}$ and $C=C_{n,a,b}$, the exact predicate is
\begin{equation}\label{eq:good}
 M>0,\qquad C>0,\qquad
 D_{n,a,b}:=4C(M-\gamma_n)-(nM)^2>0.
\end{equation}
This is the conjunction named \code{HallScalar.Good} in the source.
It is checked directly from~\eqref{eq:gamma}--\eqref{eq:curvature},
without a table of externally computed successful margins as a premise.

\begin{lemma}[Exhaustive finite criterion]\label{lem:finite}
Every triple satisfying~\eqref{eq:domain} satisfies~\eqref{eq:good}.
\end{lemma}

\begin{proof}[Verification and coverage]
Call $(a,b)$ proper if $a,b>0$ and $a+b<n$. At each of the five
dimensions, \code{Finite.lean} proves the implication from the reduced
predicate to \code{Good} for every $(a,b)$ in
$\code{Fin 15}\times\code{Fin 15}$. The reduced predicate consists
of all ordered proper pairs and the pair $(n-1,0)$.
The five declarations \code{finite\_n11} through \code{finite\_n15}
use \code{decide +kernel}. Thus Lean reduces the finite decidable
proposition using its ordinary rational arithmetic definitions.

For an arbitrary proper pair in~\eqref{eq:domain}, both coordinates
are less than $15$, so the pair lies in the checked universe. The
declaration \code{reduced\_good} proves this embedding and selects the
applicable dimension. If a coordinate is zero, both piecewise formulas
are independent of the other coordinate. Consequently their values,
and hence the predicate \code{Good}, equal those at $(n-1,0)$.
The declaration \code{all\_good} makes this split and covers every
natural pair in the domain. No transpose quotient is taken for proper
pairs, and neither marginal orientation is omitted.
\end{proof}

The finite statement is an exact computational lemma. Its proof is the
checked reduction and coverage argument just described, rather than
a printed decimal estimate for the weakest case. The complete rational
formulas and quantified checks are included in the portable source.

\begin{lemma}[Quadratic comparison]\label{lem:quadratic}
Let $C,M,g,u\in\R$ and $n\in\N$. If $C>0$ and
$4C(M-g)-(nM)^2>0$, then
\[
 Cu^2-nMu+M-g>0.
\]
\end{lemma}

\begin{proof}
The identity
\begin{align*}
 4C(Cu^2-nMu+M-g)
 &= (2Cu-nM)^2+4C(M-g)-(nM)^2
\end{align*}
has a nonnegative square and a strictly positive remainder on its
right-hand side. Division by the positive factor $4C$ proves the claim.
In Lean, \code{ring} checks the identity and ordered-field reasoning
supplies the sign conclusion.
\end{proof}

\begin{proof}[Proof of Theorem~\ref{thm:positive}]
Lemma~\ref{lem:finite}, transported from $\Q$ to $\R$, supplies
$M>0$, $C>0$, and $D_{n,a,b}>0$. Lemma~\ref{lem:quadratic} therefore
gives $Cu^2-nMu+M-\gamma_n>0$ for every real $u$.
For $u\leq 1$, Bernoulli's inequality gives
\[
 (1-u)^n\geq 1-nu.
\]
Its usual hypothesis $-u\geq-1$ follows directly from $u\leq1$.
Multiplication by $M>0$ and addition of $Cu^2-\gamma_n$ yield
\[
 Cu^2+M(1-u)^n-\gamma_n
 \geq Cu^2-nMu+M-\gamma_n>0.
\]
The formal proof invokes Mathlib's \code{one\_add\_mul\_le\_pow};
the installed version permits the weaker hypothesis $-2\leq -u$,
which also follows from the stated domain. No interval sampling,
root selection, or small-parameter assumption is used.
\end{proof}

\section{Coverage and formal artifact}

For a fixed $n$, the number of ordered proper pairs is
\[
 \sum_{s=2}^{n-1}(s-1)=\frac{(n-1)(n-2)}2.
\]
The nonfull marginal pairs are $(a,0)$ and $(0,b)$ with
$1\leq a,b\leq n-1$, giving $2(n-1)$ pairs. They share one
scalar formula. Table~\ref{tab:coverage} distinguishes these counts
from the reduced family actually checked.

\begin{table}[ht]
\centering
\begin{tabular}{@{}rrrrrr@{}}
\toprule
$n$ & Proper & Marginal & Normalized marginal & Reduced & Uncompressed\\
\midrule
11 &45&20&1&46&65\\
12 &55&22&1&56&77\\
13 &66&24&1&67&90\\
14 &78&26&1&79&104\\
15 &91&28&1&92&119\\
\midrule
Total &335&120&5&340&455\\
\bottomrule
\end{tabular}
\caption{Ordered nonfull scalar cases. ``Reduced'' adds one normalized
marginal to every proper pair. ``Uncompressed'' retains every marginal
size and orientation.}\label{tab:coverage}
\end{table}

The declarations \code{case\_counts} and \code{total\_case\_count}
prove the five reduced counts and their sum $340$ in Lean.
The combined declaration \code{finite\_certificate} includes both
universal coverage and this cardinality result. The uncompressed total
$455$ is the combinatorial calculation above, checked in the independent
source review; it is not presented as a named Lean cardinality theorem.
Universal scalar coverage follows from \code{all\_good}, regardless
of how the explanatory table is computed.

The source modules separate the obligations. \code{Definitions.lean}
contains the rational formulas and predicates. \code{Finite.lean}
proves the finite criterion and coverage. \code{Main.lean} proves
\code{HallScalar.scalar\_positive} and
\code{HallScalar.scalar\_exclusion}. Finally, \code{Audit.lean}
exposes \code{hall\_scalar\_certificate}, the same fully quantified
exclusion corollary, through the test import chain.

The fixtures in \code{Tests.lean} check the normalization
$\gamma_3=2/9$, proper and excluded cut boundaries, the normalized
marginal, and exclusion of $(0,0)$ from the reduced predicate.
They also instantiate positivity at $u=0$, $u=1$, and $u=-2$,
including a non-normalized marginal in the opposite orientation.
These fixtures supplement the universal proof; they do not replace
its coverage argument.

\section{Verification and reproducibility}

The portable package under \code{evidence/lean/} pins Lean
\code{v4.33.1} and Mathlib revision
\begin{center}
\code{0df444a360eaa60ab8c11dca51a86af692955474}.
\end{center}
Its manifest contains repository URLs and dependency revisions, without
an absolute local dependency path. With the required toolchain and
dependencies available, the package is built from that directory by
\begin{Verbatim}[fontsize=\small]
export LEAN_NUM_THREADS=1
nice -n 15 lake build
nice -n 15 lake env lean -j1 -M2048 ../FinalScalarAxiomAudit.lean
\end{Verbatim}
The configuration supplies one Lean thread and a 2048\,MB Lean memory
ceiling. A fresh environment may require the dependency cache download;
the recorded verification reused existing dependencies.
The harness also uses a lightweight axiom printer without those explicit
command-line limits; the separate direct audit above applies them.

Independent source and execution review reran the full scalar package
and the finite-only package, and directly audited the positivity theorem
as well as its corollary. Both harness records report \code{pass}.
These reviews were performed by separate assistant agents, not by
external human peer reviewers.
The four audited declarations were
\begin{Verbatim}[fontsize=\small]
HallScalar.scalar_positive
HallScalar.scalar_exclusion
HallScalar.finite_certificate
hall_scalar_certificate
\end{Verbatim}
Each printed exactly the transitive axiom list
\begin{Verbatim}[fontsize=\small]
[propext, Classical.choice, Quot.sound]
\end{Verbatim}
The accompanying \code{FinalScalarAxiomAudit.lean} checks the theorem
types and prints these closures and the imported arithmetic definitions.
The completed dependency graph contains no \code{sorryAx}, custom
axiom, or native-evaluation axiom. Rational inequalities are converted
to real inequalities using the standard cast interface; no arithmetic
operation is locally redefined.

The finite-only package uses the same definitions and finite proof as
the full package. It provides a separately executed audit of a different
target, not an independent arithmetic implementation. Independence here
refers to the reviewer and fresh executions, not to two distinct proof
algorithms. The evidence includes source hashes checked before and after
review. The review did not rebuild all Mathlib dependencies from source
or use an alternate proof-assistant kernel.

Table~\ref{tab:resources} records two observed coordinator runs from
\code{resource-measurements.json}. Both used macOS arm64, one Lean
thread, niceness 15, and the stated memory ceiling. The rows have
different incremental build states and shared machine load, so they
are observations rather than a performance comparison. They do not
measure total development effort, a clean installation, or total cost.

\begin{table}[ht]
\centering
\begin{tabular}{@{}p{0.49\linewidth}rr@{}}
\toprule
Observed build and axiom check & Wall seconds & Maximum RSS, bytes\\
\midrule
Standalone finite package, dependencies cached &36.05&1\,713\,913\,856\\
Completed real theorem, finite proofs cached &20.74&1\,534\,099\,456\\
\bottomrule
\end{tabular}
\caption{Two recorded incremental runs under shared local load.
These are not clean-build benchmarks.}\label{tab:resources}
\end{table}

\section{Related work and remaining obligations}

Exact certificate replay is an established formal-verification approach.
Siddique and Mian describe independent rational replay and a sound
Lean checker for geometric non-embeddability certificates
\cite{siddique2026machine}. That work provides a methodological
comparison for making the certificate interpretation and trusted base
explicit, rather than a runtime baseline for this finite family.
The prior Dittert formalization of Li, Xiong, and Yang uses an SOS
certificate \cite{li2026proof}; this note verifies a Hall scalar
comparison and does not reimplement their dimension-four proof.

Bernstein polynomials and de Casteljau computation have prior formal
treatments \cite{bertot2011formal,munoz2012formalization}.
Appendix~\ref{app:nearzero} separately verifies two literal expansions
in an unnormalized Bernstein basis. It uses direct polynomial identities
and positivity proofs, not a generic Bernstein checker, subdivision
theorem, or de Casteljau implementation. These additional certificates
are not used in Theorem~\ref{thm:positive}.

The boundary of the formalization is substantive. It does not construct
a Hall minorant from a matrix, prove the relevant permanent floors,
derive row or column deficit estimates, or justify stationarity and
support reductions. Marginal normalization proves equality of the
defined scalar constants; the source argument selecting a smallest row
of a matrix remains outside Lean. The full cut $(a,b)=(0,0)$ is excluded
by~\eqref{eq:domain}, and its separate analytic exclusion is not
formalized here. Nor does the artifact prove the matrix inequality's
equality characterization.

Complete coverage of the nonfull scalar domain is proved here in dimensions
11 through 15. The appendix adds only two near-zero polynomial results
in dimensions 8 and 9. Full scalar coverage in dimensions 7 through 10,
dimensions 5 and 6, and any extension above 15 remain outside
this artifact. Lower-dimensional source arguments may involve different
interfaces, including work attributed to Lu in the public proof chain;
they cannot be inferred from the scalar theorem here.\footnote{\raggedright Hongyuan
Lu's \href{https://github.com/lueluelue2006/dittert-n7-extension/tree/56397b70299278b338a286b90361180833dbf983}{\texttt{lueluelue2006/dittert-n7-extension}},
\par\noindent Source commit: \nolinkurl{56397b70299278b338a286b90361180833dbf983}.}
Formalizing additional matrix interfaces or a different certificate
family would require a separately stated result, checked coverage, and
a corresponding extension of the evidence record. The present note
claims the completed finite scalar formalization and the separately
bounded polynomial certificates, with their documented verification
boundaries only.

\appendix
\section{Separate near-zero certificates in dimensions 8 and 9}
\label{app:nearzero}

This appendix concerns only the near-zero polynomial portion of the
source marginal arguments. For $n\in\{8,9\}$, use
$\gamma_n=n!/n^n$ and define the following real constants and polynomials:
\begin{align*}
 \kappa_n&=(n-2)!\left(\frac{n-2}{(n-1)^2}\right)^{n-2},
 &L_n&=\kappa_n\left(1+\frac{1}{2n^2}\right),\\
 E_n(u)&=1-(2-\gamma_n)u,
 &H_n(u)&=(1-\gamma_n)u,\\
 T_n(u)&=E_n(u)-\frac{H_n(u)}{n-1},
 &F_n(u)&=(1-u)\left(1+\frac{u}{n-1}\right)^{n-1}.
\end{align*}
The literal polynomial checked in \code{NearZero.p} is
\begin{equation}\label{eq:nearpoly}
\begin{split}
 p_n(u)={}&\bigl(2-\gamma_n+L_n(1-u)^n-F_n(u)\bigr)E_n(u)^n\\
          &-\bigl(E_n(u)+H_n(u)\bigr)T_n(u)^{n-1}.
\end{split}
\end{equation}
There is no division by $E_n(u)$ in this definition and thus no
nonvanishing premise for $E_n$ is omitted.

\begin{theorem}[Literal near-zero positivity]\label{thm:nearzero}
For every real $u$, $0\leq u\leq 1/400$ implies $p_8(u)>0$, and
$0\leq u\leq 1/500$ implies $p_9(u)>0$.
\end{theorem}

\begin{proof}
For $(n,d)=(8,400)$ and $(9,500)$, the supplied rational coefficients
$c_{n,k}>0$ give the exact identities, for every real $x$,
\begin{equation}\label{eq:nearidentity}
 p_n(x/d)=\sum_{k=0}^{2n}c_{n,k}x^k(1-x)^{2n-k}.
\end{equation}
There are respectively 17 and 19 coefficients, with homogeneous basis
degrees 16 and 18. The binomial factors of the normalized Bernstein
basis are absorbed into $c_{n,k}$. The complete rational lists are
the literal constants in \code{Certificate8.lean} and
\code{Certificate9.lean}, under \code{evidence/near-zero/}.
For example, their initial coefficients are
\[
 c_{8,0}=\frac{77405634405}{1814205196009472},\qquad
 c_{9,0}=\frac{33970014498175}{2629465015396073472}.
\]
The declarations \code{NearZero.identity8} and \code{NearZero.identity9}
prove~\eqref{eq:nearidentity} by unfolding~\eqref{eq:nearpoly},
normalizing factorials and rational arithmetic with \code{norm\_num},
and checking the full identity with \code{ring}.

If $0\leq x<1$, all summands are nonnegative and the $k=0$
summand is strictly positive. At $x=1$ only the last summand
survives, with value $c_{n,2n}>0$. The declarations
\code{NearZero.positive8} and \code{NearZero.positive9} check these
signs exactly, including both endpoints. Finally, $0\leq u\leq1/d$
implies $0\leq du\leq1$, and $(du)/d=u$. This substitution proves
the conjunction named \code{near\_zero\_certificate} in \code{Main.lean}.
Neither the coefficient signs nor the identities are hypotheses of
that theorem.
\end{proof}

The Python generator proposes constants and source text; it is outside
the trusted proof. A separate assistant-agent review reran and audited
the final conjunction and the four identity and positivity declarations, each
with exactly \code{propext}, \code{Classical.choice}, and
\code{Quot.sound}. The portable package uses the same pinned Lean
and Mathlib versions as the scalar package. Its \code{replay.sh}
rebuilds the proof modules and runs the five-declaration axiom audit;
its README gives the standalone command and dependency requirements.

These certificates do not prove the remaining support intervals,
their monotone tails, support compression, matrix reductions, or an
$n=7$ case. In particular, they are neither full dimension-8 or
dimension-9 scalar coverage nor full Dittert proofs. Their contribution
is the reproducible formal verification of these two existing,
explicitly bounded polynomial inequalities.

\bibliographystyle{plainnat}
\bibliography{draft/references}
\end{document}
